\documentclass[11pt]{article}
\usepackage[T2A]{fontenc}
\usepackage[utf8]{inputenc}
\usepackage[ukrainian,english]{babel}
\usepackage[margin=1in]{geometry}
\usepackage{amsmath,amssymb}
\usepackage{enumitem}
\usepackage{hyperref}
\usepackage{tikz}
\usepackage{microtype}
\usepackage{longtable}
\usepackage{booktabs}

\setlength{\parindent}{0pt}
\setlength{\parskip}{0.6\baselineskip}
\setlength{\emergencystretch}{3em}
\hbadness=10000
\sloppy

\title{Багатовимірна матриця складності для оцінювання великих мовних моделей у задачах символьної регресії}
\author{}
\date{}

\begin{document}

\maketitle

\section*{Анотація}

Швидкий прогрес великих мовних моделей (LLM) та поява архітектур, орієнтованих на міркування (reasoning-first, тобто з пріоритетом міркування), актуалізують питання їх застосування у задачах наукового пошуку, зокрема у символьній регресії --- відновленні аналітичної форми математичних залежностей за емпіричними даними. Проте валідність існуючих бенчмарків є обмеженою через статичність тестових наборів та високу ймовірність присутності розв'язків у навчальних корпусах моделей, що призводить до хибної оцінки здатності моделей до узагальнення. У статті обґрунтовано необхідність переходу від статичних тестів до процедурної генерації завдань та запропоновано нову теоретичну методологію оцінювання --- багатовимірну таксономію математичних виразів. Наукова новизна роботи полягає в розробці Матриці складності (Complexity Matrix), що класифікує функції за трьома ортогональними вимірами: структурна складність дерева виразів, семантичний простір операторів та топологічні особливості поведінки функції. Вперше запропоновано використовувати топологічні властивості (асимптотика, розриви, періодичність) як окремий критерій для оцінки глибини розуміння моделлю математичної природи даних. Запропонована концепція створює теоретичне підґрунтя для розробки адаптивних систем оцінювання, здатних диференціювати ефект запам'ятовування від реальних навичок математичного міркування у сучасних LLM.

\textbf{Ключові слова:} символьна регресія, великі мовні моделі, генералізація, таксономія функцій, обчислювальні графи, математичне міркування, оцінювання ШІ.

\section*{Вступ}

Останні роки ознаменувалися фундаментальним зсувом у галузі штучного інтелекту: великі мовні моделі (LLM) еволюціонували від імовірнісних генераторів тексту до інструментів, здатних до складних ланцюжків міркувань (міркування у форматі Chain-of-Thought, CoT). Поява моделей класу reasoning-first (із пріоритетом міркування), таких як серія OpenAI o1 та DeepSeek-R1, відкрила нові горизонти для застосування ШІ у фундаментальних науках, формуючи парадигму AI for Science~[1]. Одним із найбільш перспективних напрямів у цьому контексті є символьна регресія (Symbolic Regression, SR) --- задача пошуку точного аналітичного виразу, що описує набір експериментальних даних~[2]. На відміну від традиційних нейронних мереж, які діють за принципом <<чорної скриньки>>, символьна регресія забезпечує інтерпретованість результату, що є критично важливим для виявлення нових фізичних законів~[3]. Дослідження 2024--2025 років демонструють, що сучасні LLM здатні виступати як наскрізні (end-to-end) генератори наукових гіпотез, використовуючи апріорні знання про математичні структури, що перевершує ефективність класичних методів генетичного програмування у певних класах задач~[4,5].

Проте, попри оптимістичні результати, наукова спільнота стикається з кризою оцінювання (evaluation crisis). Валідність існуючих методик бенчмаркінгу в математичних задачах викликає сумніви через проблему <<забруднення даних>> (data contamination). Традиційні набори даних, такі як Feynman Symbolic Regression Dataset або SRBench, є відкритими та статичними. Враховуючи масштаби навчальних корпусів сучасних моделей, існує висока ймовірність, що ці формули вже були опрацьовані моделями під час попереднього навчання. Це підтверджується дослідженнями кінця 2024 року, зокрема роботою Mirzadeh et al.
під назвою GSM-Symbolic~[9], яка показує, що навіть передові моделі часто покладаються на ймовірнісне зіставлення шаблонів, а не на логічне виведення. У контексті символьної регресії це означає, що високі показники точності апроксимації (наприклад, $R^2>0.99$) можуть бути наслідком відтворення запам'ятованого розв'язку, а не успішної генералізації.

Наукова проблема ускладнюється обмеженістю існуючих метрик. Більшість підходів фокусуються на чисельній точності (MSE) або простоті виразу, ігноруючи структурну семантику функції. Модель може апроксимувати складну топологічну особливість, наприклад полюс, поліномом високого ступеня, формально досягаючи високої точності на обмеженому інтервалі, але фундаментально помиляючись у фізичній суті процесу~[8,10]. Наразі відсутня систематизована таксономія, яка класифікувала б математичні вирази не лише за їх синтаксичною довжиною, а й за топологічними властивостями та когнітивною складністю реконструкції.

Метою даної роботи є розробка теоретичної методології для створення стійких до витоку даних систем оцінювання. Для досягнення цієї мети запропоновано концепцію Багатовимірної матриці складності (Multidimensional Complexity Matrix), що передбачає класифікацію функцій у трьох ортогональних вимірах: структурна складність графа, семантичний простір операторів та топологічні ознаки поведінки. Введення останнього виміру є ключовою науковою новизною, що дозволяє оцінювати здатність LLM розпізнавати якісну поведінку функції. Така таксономія створює фундамент для процедурної генерації позарозподільних (out-of-distribution) завдань, що унеможливлює ефект механічного запам'ятовування та дозволяє об'єктивно оцінити навички математичного міркування моделей.

Принципово важливо, що дана стаття має концептуально-теоретичний характер. У межах цієї публікації не подається програмна реалізація генератора, не виконується масштабне експериментальне порівняння моделей та не заявляються емпіричні висновки щодо переваг конкретних LLM. Ці етапи розглядаються як окремий наступний цикл досліджень.

\section{Багатовимірна матриця складності}

Центральним елементом дослідження є розробка методології класифікації завдань символьної регресії, яка відходить від одновимірних метрик (таких як довжина виразу) на користь векторного представлення складності. Вводиться поняття Багатовимірної матриці складності (Multidimensional Complexity Matrix, MCM), що розглядає математичний вираз як об'єкт із трьома ортогональними властивостями: структурною топологією графа обчислень, семантичною ентропією операторів та якісною поведінкою функції.

\subsection{Формалізація простору пошуку}

Нехай $\mathcal{F}$ --- простір аналітичних функцій $f:\mathbb{R}^n\to\mathbb{R}$, що підлягають відновленню. Кожна функція $f\in\mathcal{F}$ може бути представлена у вигляді синтаксичного дерева (expression tree)
\[
T_f=(V,E),
\]
де $V$ --- множина вершин (оператори та операнди), а $E$ --- множина ребер, що визначають порядок обчислень.

Традиційний підхід до оцінки складності базується на скалярній метриці
\[
L(f)=|V|,
\]
що відображає лише кількість вузлів. У даній роботі обґрунтовано необхідність використання відображення
\[
\Psi:\mathcal{F}\to\mathbb{R}^3,
\]
яке ставить у відповідність кожній функції вектор складності
\[
\mathbf{C}_f=\langle \alpha,\beta,\gamma\rangle,
\]
де $\alpha$ --- показник структурної складності (вісь A), $\beta$ --- семантичний індекс (вісь B), $\gamma$ --- топологічний дескриптор (вісь C).

\subsection{Вісь A: структурна складність (Structural Complexity, \texorpdfstring{$\alpha$}{alpha})}

Цей вимір характеризує топологію графа $T_f$ як абстрактної обчислювальної структури, незалежно від математичного змісту вузлів. Показник визначає навантаження на контекстне вікно моделі та необхідну глибину логічного виведення.

Для класифікації виділено чотири дискретні рівні:

\begin{itemize}[leftmargin=1.5em]
  \item \textbf{$A_0$ (атомарні та тривіальні структури).} Графи мінімальної глибини ($d\leq 1$): ідентичність $f(x)=x$, константи, прості унарні операції. Цей рівень використовується як контрольна група (базовий рівень) для виявлення ефекту надмірного ускладнення міркування (overthinking) та переускладнення моделі (overfitting) у моделях із пріоритетом міркування.
  \item \textbf{$A_1$ (поліноміальні та адитивні структури).} Графи з великою шириною при помірній глибині: лінійні комбінації базисних функцій або поліноми.
  \item \textbf{$A_2$ (композиційні структури).} Вкладені функції, що утворюють ланцюги послідовних залежностей, наприклад $f(g(h(x)))$. Когнітивне навантаження високе, оскільки потрібно формувати багатоступеневі проміжні кроки міркування.
  \item \textbf{$A_3$ (рекурсивні та фрактальні структури).} Вирази з аномальною глибиною ($d>5$) або специфічною вкладеною структурою (ланцюгові дроби, вкладені радикали). Рівень тестує здатність моделі обробляти довгі залежності та рекурсію.
\end{itemize}

\subsection{Вісь B: семантичний простір (Semantic Operator Set, \texorpdfstring{$\beta$}{beta})}

Вимір визначає алфавіт операторів $\Omega$, доступних для побудови функції. Складність зростає зі збільшенням ентропії простору пошуку та зменшенням частоти зустрічальності операторів у навчальних корпусах.

\begin{itemize}[leftmargin=1.5em]
  \item \textbf{$B_0$ (арифметичний базис).} $\Omega_0=\{+, -, \times, \div\}$. Задачі цього рівня перевіряють базову числову грамотність.
  \item \textbf{$B_1$ (гладкі трансцендентні функції).} $\Omega_1=\Omega_0\cup\{\sin,\cos,\exp,\log,\mathrm{pow}\}$. Це стандартний набір для опису гладких фізичних процесів.
  \item \textbf{$B_2$ (недиференційовні та логічні оператори).} $\Omega_2=\Omega_1\cup\{|x|,\mathrm{sgn}(x),\lfloor x\rfloor,\max(x,y)\}$. Введення цих операторів формує негладкість і вимагає умовного логічного розгалуження.
  \item \textbf{$B_3$ (спеціальні функції).} $\Omega_3=\Omega_2\cup\{\Gamma(x),J_\nu(x),\mathrm{erf}(x)\}$. Виділення цього класу дозволяє ізольовано оцінити предметну експертизу та здатність моделі працювати з рідкісним математичним словником.
\end{itemize}

\subsection{Вісь C: топологічні ознаки (Topological Feature Salience, \texorpdfstring{$\gamma$}{gamma})}

Даний вимір є ключовою науковою новизною дослідження. Пропонується оцінювати не лише синтаксис формули, а й топологічний портрет функції у просторі ознак. Гіпотеза полягає в тому, що справжнє математичне розуміння передбачає здатність детектувати якісні характеристики поведінки за числовими даними.

\begin{itemize}[leftmargin=1.5em]
  \item \textbf{$C_0$ (регулярність).} Функції гладкі, монотонні або мають просту опуклість; локальна апроксимація рядом Тейлора працює ефективно.
  \item \textbf{$C_1$ (інваріантність та симетрія).} Функції демонструють глобальні патерни повторюваності: періодичність $f(x+T)=f(x)$, парність або непарність.
  \item \textbf{$C_2$ (асимптотична поведінка).} Наявні ізольовані сингулярності з передбачуваною локальною поведінкою: полюси, горизонтальні асимптоти або вертикальні дотичні. Рівень тестує розуміння меж і нескінченності.
  \item \textbf{$C_3$ (сингулярності та хаотична поведінка).} Функції мають якісні переходи та акумуляцію особливих точок: розриви, cusps, нескінченні осциляції на скінченному інтервалі (наприклад, $\sin(1/x)$ при $x\to 0$), граничні множини полюсів.
\end{itemize}

На відміну від $C_2$, де достатньо детектувати окрему сингулярність, клас $C_3$ вимагає розуміння глобальної структури множини особливих точок та їх взаємодії.

\subsection{Операціоналізація класифікації (Axes A/B/C)}

Для забезпечення відтворюваності таксономії вводиться детерміноване відображення
\[
K:\mathcal{F}\to\{0,1,2,3\}^3,\qquad K(f)=\langle A(f),B(f),C(f)\rangle.
\]

\textbf{Операціоналізація осі A.}
Для дерева виразу $T_f=(V,E)$ визначаються ознаки
\[
d(f)=\mathrm{depth}(T_f),\qquad n(f)=|V|,\qquad r(f)\in\{0,1\},
\]
де $r(f)=1$ означає наявність рекурсивних або самоподібних шаблонів (ланцюгові дроби, вкладені радикали, повторна композиція одного типу глибини $\ge 2$). Класифікація задається правилом
\[
A(f)=
\begin{cases}
0,& d(f)\le 1 \land n(f)\le 3,\\
1,& 1<d(f)\le 3 \land r(f)=0,\\
2,& 3<d(f)\le 5 \land r(f)=0,\\
3,& d(f)>5 \lor r(f)=1.
\end{cases}
\]

\textbf{Операціоналізація осі B.}
Нехай $\Omega_f$ --- множина операторів, що фактично присутні у виразі $f$. Вводиться рівнева функція
\[
\lambda(op)=
\begin{cases}
0,& op\in\Omega_0,\\
1,& op\in\Omega_1\setminus\Omega_0,\\
2,& op\in\Omega_2\setminus\Omega_1,\\
3,& op\in\Omega_3\setminus\Omega_2,
\end{cases}
\]
а семантичний клас визначається як
\[
B(f)=\max_{op\in\Omega_f}\lambda(op).
\]

\textbf{Операціоналізація осі C.}
Визначаються індикатори: $S(f)$ --- множина особливих точок; $p(f)$ --- періодичність ($1$, якщо $\exists T>0: f(x+T)=f(x)$); $\sigma(f)$ --- симетрія ($1$, якщо $f(-x)=f(x)$ або $f(-x)=-f(x)$); $o(f)$ --- наявність нескінченних осциляцій у скінченному околі; $a(f)$ --- наявність точки акумуляції множини $S(f)$ у скінченному інтервалі. Тоді
\[
C(f)=
\begin{cases}
0,& S(f)=\varnothing \land p(f)=0 \land o(f)=0,\\
1,& S(f)=\varnothing \land (p(f)=1 \lor \sigma(f)=1),\\
2,& 0<|S(f)|<\infty \land a(f)=0 \land o(f)=0,\\
3,& a(f)=1 \lor o(f)=1 \lor (\exists I\subset\mathbb{R}_{\text{bounded}}:|S(f)\cap I|=\infty).
\end{cases}
\]

У випадку прикордонних або змішаних ознак застосовується консервативне правило розв'язання конфлікту: функція відноситься до найвищого сумісного рівня за відповідною віссю.

Для повноти класифікації нижче наведено ілюстративну бібліотеку прикладів для всіх 64 векторів $\langle A_i,B_j,C_k\rangle$.

{\small
\setlength{\tabcolsep}{6pt}
\renewcommand{\arraystretch}{1.12}
\newcommand{\mcmrow}[2]{#1 & #2\\[-0.15em]\cmidrule(lr){1-2}\addlinespace[0.15em]}
\begin{longtable}{p{0.22\textwidth}p{0.72\textwidth}}
\caption{Ілюстративні приклади функцій для повного простору $4\times 4\times 4$.}\\
\toprule
\textbf{Вектор складності} & \textbf{Приклад функції $f(x)$}\\
\midrule
\endfirsthead
\toprule
\textbf{Вектор складності} & \textbf{Приклад функції $f(x)$}\\
\midrule
\endhead
\midrule
\multicolumn{2}{r}{\textit{Продовження на наступній сторінці}}\\
\midrule
\endfoot
\bottomrule
\endlastfoot

\multicolumn{2}{l}{\textit{Блок $C_0$ (регулярні функції)}}\\
\mcmrow{$\langle A_0,B_0,C_0\rangle$}{$x+1$}
\mcmrow{$\langle A_1,B_0,C_0\rangle$}{$x(x+1)+1$}
\mcmrow{$\langle A_2,B_0,C_0\rangle$}{$\dfrac{(x+1)(x+2)}{1+x(x+1)}$}
\mcmrow{$\langle A_3,B_0,C_0\rangle$}{$\dfrac{1}{1+\dfrac{1}{1+\dfrac{1}{1+x(x+1)}}}$}
\mcmrow{$\langle A_0,B_1,C_0\rangle$}{$e^x$}
\mcmrow{$\langle A_1,B_1,C_0\rangle$}{$e^x+x$}
\mcmrow{$\langle A_2,B_1,C_0\rangle$}{$\sin\!\left(e^{x+1}\right)+x$}
\mcmrow{$\langle A_3,B_1,C_0\rangle$}{$\log\!\big(1+e^{\sin(\log(1+x^2))}\big)+x$}
\mcmrow{$\langle A_0,B_2,C_0\rangle$}{$\max(x,x+1)$}
\mcmrow{$\langle A_1,B_2,C_0\rangle$}{$|x|^2+x$}
\mcmrow{$\langle A_2,B_2,C_0\rangle$}{$\log\!\big(1+e^{|x|^2+x}\big)$}
\mcmrow{$\langle A_3,B_2,C_0\rangle$}{$\log\!\big(1+e^{\sin(\max(x,x+1))}\big)+x$}
\mcmrow{$\langle A_0,B_3,C_0\rangle$}{$\mathrm{erf}(x+1)$}
\mcmrow{$\langle A_1,B_3,C_0\rangle$}{$\mathrm{erf}(x+1)+x$}
\mcmrow{$\langle A_2,B_3,C_0\rangle$}{$\mathrm{erf}(\sin(x+1))+x$}
\mcmrow{$\langle A_3,B_3,C_0\rangle$}{$\mathrm{erf}\!\big(\log(1+e^{\sin(\log(1+x^2))})\big)+x$}
\midrule
\multicolumn{2}{l}{\textit{Блок $C_1$ (симетрія/інваріантність)}}\\
\mcmrow{$\langle A_0,B_0,C_1\rangle$}{$x-x$}
\mcmrow{$\langle A_1,B_0,C_1\rangle$}{$x^2+1$}
\mcmrow{$\langle A_2,B_0,C_1\rangle$}{$\dfrac{x^2+1}{1+\dfrac{1}{x^2+2}}$}
\mcmrow{$\langle A_3,B_0,C_1\rangle$}{$\dfrac{1}{1+\dfrac{1}{1+\dfrac{1}{x^2+1}}}$}
\mcmrow{$\langle A_0,B_1,C_1\rangle$}{$\sin x$}
\mcmrow{$\langle A_1,B_1,C_1\rangle$}{$\sin x+\cos x$}
\mcmrow{$\langle A_2,B_1,C_1\rangle$}{$\sin\!\left(\cos\!\left(e^{-x^2}\right)\right)$}
\mcmrow{$\langle A_3,B_1,C_1\rangle$}{$\sin\!\big(\cos(\sin(\cos(\sin(\cos x))))\big)$}
\mcmrow{$\langle A_0,B_2,C_1\rangle$}{$\max(x,x)$}
\mcmrow{$\langle A_1,B_2,C_1\rangle$}{$|x|^2$}
\mcmrow{$\langle A_2,B_2,C_1\rangle$}{$\sin\!\left(\cos\!\left(e^{-|x|^2}\right)\right)$}
\mcmrow{$\langle A_3,B_2,C_1\rangle$}{$\sin\!\big(\cos(\sin(\cos(\sin(\max(x,x)))))\big)$}
\mcmrow{$\langle A_0,B_3,C_1\rangle$}{$\mathrm{erf}(x)$}
\mcmrow{$\langle A_1,B_3,C_1\rangle$}{$J_0(x)$}
\mcmrow{$\langle A_2,B_3,C_1\rangle$}{$\mathrm{erf}(\sin x)$}
\mcmrow{$\langle A_3,B_3,C_1\rangle$}{$\mathrm{erf}\!\big(\sin(\cos(\sin(\cos(\sin x))))\big)$}
\midrule
\multicolumn{2}{l}{\textit{Блок $C_2$ (ізольовані сингулярності)}}\\
\mcmrow{$\langle A_0,B_0,C_2\rangle$}{$\dfrac{1}{x}$}
\mcmrow{$\langle A_1,B_0,C_2\rangle$}{$\dfrac{1}{x-1}$}
\mcmrow{$\langle A_2,B_0,C_2\rangle$}{$\dfrac{x+1}{1+\dfrac{1}{x-1}}$}
\mcmrow{$\langle A_3,B_0,C_2\rangle$}{$\dfrac{1}{1+\dfrac{1}{1+\dfrac{1}{x-1}}}$}
\mcmrow{$\langle A_0,B_1,C_2\rangle$}{$\log x$}
\mcmrow{$\langle A_1,B_1,C_2\rangle$}{$\dfrac{\sin x}{x-1}$}
\mcmrow{$\langle A_2,B_1,C_2\rangle$}{$\dfrac{e^{\sin x}}{1+\dfrac{1}{x-1}}$}
\mcmrow{$\langle A_3,B_1,C_2\rangle$}{$\dfrac{\log\!\big(1+e^{\sin(\cos x)}\big)}{1+\dfrac{1}{1+\dfrac{1}{x-1}}}$}
\mcmrow{$\langle A_0,B_2,C_2\rangle$}{$\dfrac{1}{|x|}$}
\mcmrow{$\langle A_1,B_2,C_2\rangle$}{$\dfrac{|x|^2+1}{x-1}$}
\mcmrow{$\langle A_2,B_2,C_2\rangle$}{$\dfrac{\log(1+|x|)}{1+\dfrac{1}{x-1}}$}
\mcmrow{$\langle A_3,B_2,C_2\rangle$}{$\dfrac{\max(x,x+1)}{1+\dfrac{1}{1+\dfrac{1}{x-1}}}$}
\mcmrow{$\langle A_0,B_3,C_2\rangle$}{$\Gamma(x)$}
\mcmrow{$\langle A_1,B_3,C_2\rangle$}{$\Gamma(x)+e^{-x}$}
\mcmrow{$\langle A_2,B_3,C_2\rangle$}{$\dfrac{\Gamma(x+2)}{1+\dfrac{1}{x-1}}$}
\mcmrow{$\langle A_3,B_3,C_2\rangle$}{$\dfrac{\mathrm{erf}(\Gamma(x+2))}{1+\dfrac{1}{1+\dfrac{1}{x-1}}}$}
\midrule
\multicolumn{2}{l}{\textit{Блок $C_3$ (розриви/акумуляція особливостей)}}\\
\mcmrow{$\langle A_0,B_0,C_3\rangle$}{\textit{Н/З: клас порожній для }$\Omega_0$}
\mcmrow{$\langle A_1,B_0,C_3\rangle$}{\textit{Н/З: клас порожній для }$\Omega_0$}
\mcmrow{$\langle A_2,B_0,C_3\rangle$}{\textit{Н/З: клас порожній для }$\Omega_0$}
\mcmrow{$\langle A_3,B_0,C_3\rangle$}{\textit{Н/З: клас порожній для }$\Omega_0$}
\mcmrow{$\langle A_0,B_1,C_3\rangle$}{$\sin\!\left(\dfrac{1}{x}\right)$}
\mcmrow{$\langle A_1,B_1,C_3\rangle$}{$\dfrac{\sin(1/x)}{x-1}$}
\mcmrow{$\langle A_2,B_1,C_3\rangle$}{$e^{\sin(1/x)}+\log(1+x^2)$}
\mcmrow{$\langle A_3,B_1,C_3\rangle$}{$\dfrac{\sin\!\left(\dfrac{1}{x-\sin x}\right)+e^{-x}}{x-\sin x}$}
\mcmrow{$\langle A_0,B_2,C_3\rangle$}{$|x|$}
\mcmrow{$\langle A_1,B_2,C_3\rangle$}{$\max(x,\lfloor x\rfloor)$}
\mcmrow{$\langle A_2,B_2,C_3\rangle$}{$e^{\sin(\lfloor x\rfloor)}+\dfrac{1}{x}$}
\mcmrow{$\langle A_3,B_2,C_3\rangle$}{$\left\lfloor\dfrac{1}{x-\sin(1/x)}\right\rfloor$}
\mcmrow{$\langle A_0,B_3,C_3\rangle$}{$\Gamma(1/x)$}
\mcmrow{$\langle A_1,B_3,C_3\rangle$}{$\mathrm{erf}(\sin(1/x))$}
\mcmrow{$\langle A_2,B_3,C_3\rangle$}{$\Gamma(\sin(1/x)+2)$}
\mcmrow{$\langle A_3,B_3,C_3\rangle$}{$\dfrac{\mathrm{erf}(\Gamma(1/x+2))}{x-\sin x}$}
\end{longtable}
}

За заданою граматикою операторів блок $B_0\times C_3$ є порожнім: для скінченного виразу на $\Omega_0=\{+, -, \times, \div\}$ не формується поведінка класу $C_3$ (акумуляція сингулярностей або нескінченні осциляції на скінченному інтервалі).

\subsection{Інтеграція вимірів}

Запропонована таксономія формує простір із $4\times4\times4=64$ унікальних класів складності. Кожен вектор
\[
\langle A_i,B_j,C_k\rangle
\]
використовується як координата для стрес-тестування конкретної когнітивної здібності LLM.

\begin{figure}[h]
\centering
\begin{tikzpicture}[scale=0.9]
  \draw[thick] (0,0) rectangle (4,4);
  \draw[thick] (1.2,1.2) rectangle (5.2,5.2);
  \draw[thick] (0,0) -- (1.2,1.2);
  \draw[thick] (4,0) -- (5.2,1.2);
  \draw[thick] (0,4) -- (1.2,5.2);
  \draw[thick] (4,4) -- (5.2,5.2);

  \foreach \x in {1,2,3} {
    \draw[gray!70] (\x,0) -- (\x,4);
    \draw[gray!70] (\x+1.2,1.2) -- (\x+1.2,5.2);
  }
  \foreach \y in {1,2,3} {
    \draw[gray!70] (0,\y) -- (4,\y);
    \draw[gray!70] (1.2,\y+1.2) -- (5.2,\y+1.2);
  }

  \fill[red] (0.5,0.5) circle (0.09);
  \fill[magenta] (4.7,3.7) circle (0.09);

  \node[below] at (2,-0.25) {Вісь A};
  \node[left] at (-0.25,2) {Вісь B};
  \node[right] at (5.45,5.0) {Вісь C};
  \node[red,anchor=west] at (5.45,0.8) {$\langle A_0,B_0,C_0\rangle$ (базовий рівень)};
  \node[magenta,anchor=west] at (5.45,0.3) {$\langle A_3,B_1,C_3\rangle$ (показовий приклад)};
\end{tikzpicture}
\caption{Схематична візуалізація простору MCM ($4\times4\times4$).}
\end{figure}

Варто зазначити, що тривимірна класифікація не вичерпує усіх аспектів складності символьної регресії. Перспективним додатковим виміром є параметрична чутливість --- ступінь залежності форми функції від точних значень числових констант. Наприклад, $\sin(100x)$ та $\sin(x)$ належать до одного класу $\langle A,B,C\rangle$, але ідентифікація частотного параметра у першому випадку суттєво складніша. У подальших дослідженнях цей аспект може бути формалізований як четверта вісь (вісь D).

\section{Методологія процедурної генерації та оцінювання}

Запропонована таксономія є фундаментом для побудови алгоритмічного конвеєра генерації тестових завдань. Мета полягає у створенні синтетичних наборів даних, що гарантують позарозподільний (out-of-distribution) характер прикладів і забезпечують контроль над параметрами складності. У межах цієї статті конвеєр задається на рівні теоретичної специфікації, без програмної реалізації.

У межах даної методології розглядається реконструкція функції на безшумних даних (noiseless data). Вплив стохастичного шуму свідомо винесено за межі цієї роботи, оскільки поточний етап спрямовано на ізоляцію фундаментальної здатності моделей до математичного міркування та розпізнавання топології. Розширення на шумові режими розглядається як наступний етап дослідницької програми.

\subsection{Алгоритм генерації <<Зворотний потік>> (Reverse-Flow Generation)}

На відміну від емпіричного підходу <<дані $\rightarrow$ апроксимація>>, застосовується генерація <<формула $\rightarrow$ дані>>. Процес визначається цільовим вектором $\langle A_i,B_j,C_k\rangle$ і на концептуальному рівні поділяється на три етапи.

\textbf{Етап 1: синтаксичний синтез (Axes A \& B).}
Визначається процедура побудови множини кандидатних дерев виразів за допомогою стохастичного зростання дерев з обмеженнями на глибину та дозволені оператори. Для забезпечення валідності високих рівнів складності передбачається примусове включення вузлів, які відповідають заданому класу.

\textbf{Етап 2: топологічна фільтрація (вісь C).}
Передбачається верифікація кандидатів із використанням систем комп'ютерної алгебри (CAS)~[13] і чисельного аналізу: детекція сингулярностей, аналіз асимптотики, перевірка періодичності.

Розглянемо приклад кандидата
\[
f_{\mathrm{cand}}(x)=\frac{\sin\!\left(\frac{1}{x^2+\cos(1/x)}\right)+e^{-x}}{x-\sin x}.
\]
Для нього структурний аналіз виявляє глибоку вкладеність (клас $A_3$), семантичний аналіз фіксує стандартні трансцендентні оператори (клас $B_1$), а топологічний аналіз виявляє сингулярності й нескінченні осциляції (клас $C_3$). Отже, кандидат класифікується як $\langle A_3,B_1,C_3\rangle$.

\textbf{Етап 3: генерація даних.}
Для затвердженої функції генерується множина точок $(x_i,y_i)$ у заданому діапазоні. Значення $y_i=f(x_i)$ обчислюються з підвищеною точністю (FP64/FP128). Використання безшумних даних дозволяє інтерпретувати помилку моделі як помилку структурної ідентифікації.

\subsection{Протокол оцінювання}

Традиційної метрики $R^2$ недостатньо, оскільки вона не розрізняє <<правильний фізичний закон>> та <<локально точну поліноміальну апроксимацію>>. Пропонується ієрархічний протокол:

\textbf{1) Точне символьне співпадіння (Exact Symbolic Match, ESM).}
Бінарна метрика:
\[
\mathrm{ESM}(\hat{f},f)=\mathbb{I}\big[\hat{f}-f\equiv 0\big].
\]
Проблема еквівалентності виразів у загальному випадку алгоритмічно нерозв'язна (теорема Річардсона)~[14], тому практична реалізація ESM базується на комбінації алгебраїчного спрощення в CAS~[13] і ймовірнісної чисельної верифікації на великій множині випадкових точок.

Для відтворюваності доцільно фіксувати процедуру ESM-перевірки: $N=10^4$ випадкових точок на об'єднанні інтервалів без $\varepsilon$-околів сингулярностей ($\varepsilon=10^{-6}$), з критерієм прийняття $\max_i|\hat f(x_i)-f(x_i)|<10^{-8}$.

\textbf{2) Структурна подібність (Tree Edit Distance, TED).}
Якщо $\mathrm{ESM}=0$, відстань редагування дерев кількісно оцінює близькість гіпотези до істини (помилка у константі, знаку або окремому операторі).

\textbf{3) Екстраполяційна стійкість (Extrapolation Consistency).}
Оцінюється MSE на діапазоні, що лежить за межами інтервалу навчальних точок. Низька помилка екстраполяції свідчить про коректне відтворення топології функції.

\subsection{Калібрування та контроль (базовий контроль)}

Методологія передбачає обов'язкове використання контрольної групи задач класу $\langle A_0,B_0,C_0\rangle$ (тривіальні функції типу $y=x$ або $y=x^2$). Відсутність 100\% точності на цьому рівні свідчить про проблеми калібрування моделі та її схильність до галюцинування складності.

\section{Дискусія}

Запропонована методологія MCM створює передумови для переходу від агрегованого оцінювання точності до якісного профілювання когнітивних здібностей LLM у математичному домені.

\subsection{Від запам'ятовування до абдуктивного міркування}

Фундаментальною проблемою застосування LLM у науці є непрозорість механізму формування відповіді. Висока результативність на статичних тестах часто корелює з обсягом навчальних даних, а не з глибиною розуміння. Успіх у складних топологічних класах $C_2$--$C_3$ може інтерпретуватися як ознака абдуктивного міркування --- здатності формулювати гіпотези, що пояснюють спостережувані аномалії.

\subsection{Діагностика когнітивних викривлень}

Поєднання контрольної групи тривіальних функцій із задачами високої складності дозволяє виявляти специфічні для LLM когнітивні викривлення, зокрема схильність до надмірного ускладнення міркування (overthinking). Для кількісного аналізу цього ефекту доцільно використовувати метрику структурної надлишковості (Structural Redundancy Score, SRS).

\subsection{Обмеження підходу}

По-перше, фокус на безшумному режимі (noiseless) ідеалізує задачу; реальні експерименти вимагають урахування шуму та статистичної невизначеності. По-друге, еквівалентність аналітичних виразів у загальному випадку нерозв'язна~[14], тому практичний ESM-аналіз є наближеним, хоча й достатньо надійним для більшості класів у межах запропонованої таксономії.

Наукова новизна підходу полягає не в окремих компонентах (синтаксична глибина, класи гладкості), а в їх інтеграції у єдину метричну систему, адаптовану саме для оцінювання когнітивних можливостей LLM у задачах символьної регресії.

\section*{Висновки}

У роботі розглянуто проблему валідності оцінювання великих мовних моделей у задачах символьної регресії та показано обмеження статичних бенчмарків в умовах забруднення даних (data contamination).

Основним науковим результатом є теоретичне обґрунтування Багатовимірної матриці складності, що класифікує задачі за трьома ортогональними осями: структурна складність графа обчислень (вісь A), семантичний простір операторів (вісь B) і топологічні ознаки поведінки функції (вісь C).

Запропонована методологія процедурної генерації завдань <<Зворотний потік>> (Reverse-Flow Generation) та ієрархічний протокол оцінювання (ESM, TED, Extrapolation Consistency) формують фундамент для побудови адаптивних систем тестування нового покоління.

Перспективи подальших досліджень включають програмну реалізацію генератора, побудову динамічного бенчмарку на основі MCM, а також масштабне порівняння провідних LLM (GPT, Claude, Llama) у просторі складності $\langle A,B,C\rangle$.

\section*{Список використаних джерел}

\begin{enumerate}[label={[{\arabic*}]}]
  \item OpenAI. \textit{OpenAI o1 System Card}. OpenAI Technical Reports, 2024. \url{https://openai.com}.
  \item Schmidt, M., \& Lipson, H. Distilling free-form natural laws from experimental data. \textit{Science}, 324(5923), 81--85, 2009.
  \item Udrescu, S. M., \& Tegmark, M. AI Feynman: A physics-inspired method for symbolic regression. \textit{Science Advances}, 6(16), eaay2631, 2020.
  \item Shojaee, P., et al. LLM-SR: Scientific Equation Discovery via Programming with Large Language Models. \textit{arXiv preprint}, 2024.
  \item Merler, M., et al. In-Context Symbolic Regression: Leveraging Large Language Models for Function Discovery. \textit{Proceedings of ACL}, 2024.
  \item Golchin, S., \& Surdeanu, M. Data Contamination Quiz: A Tool to Detect and Estimate Contamination in Large Language Models. \textit{arXiv preprint}, 2024.
  \item Udrescu, S. M., et al. AI Feynman 2.0: Pareto-optimal symbolic regression exploiting graph modularity. \textit{NeurIPS}, 33, 4860--4871, 2020.
  \item La Cava, W., et al. Contemporary symbolic regression methods and their relative performance. \textit{NeurIPS Datasets and Benchmarks Track}, 2021.
  \item Mirzadeh, I., et al. GSM-Symbolic: Understanding the Limitations of Mathematical Reasoning in Large Language Models. \textit{arXiv preprint}, 2024.
  \item Yang, Y., et al. Memorization vs. Reasoning in Large Language Models: A Comprehensive Survey. \textit{Transactions on Machine Learning Research}, 7(1), 112--145, 2025.
  \item Koza, J. R. \textit{Genetic Programming: On the Programming of Computers by Means of Natural Selection}. MIT Press, 1992.
  \item Petersen, B. K., et al. Deep symbolic regression: Recovering mathematical expressions from data via risk-seeking policy gradients. \textit{ICLR}, 2021.
  \item Meurer, A., et al. SymPy: symbolic computing in Python. \textit{PeerJ Computer Science}, 3, e103, 2017.
  \item Richardson, D. Some undecidable problems involving elementary functions of a real variable. \textit{Journal of Symbolic Logic}, 33(4), 514--520, 1968.
\end{enumerate}

\end{document}